\section{Lecture 8: Huber Pairs and Analytic Rings}
\label{lec:8}

Scholze's last lecture (before the course turns to geometry proper) sets up the
general notion of an \emph{analytic ring} and compares it with Huber's theory of
adic spaces. The material has two halves. First, the abstract definition: an
analytic ring is an underlying light condensed ring $\ur A$ together with a
notion of ``complete'' modules --- a full abelian subcategory
$\Mod(A)\subseteq\Cond(\ur A)$ closed under enough operations to behave like the
solid or $\mathbb{Z}[T]$-solid categories of \cref{lec:5,lec:7}. One derives from
this a base-change left adjoint, a symmetric monoidal structure, a derived
category $D(A)$ with a $t$-structure, and a derived symmetric monoidal tensor
product. Second, the comparison with Huber: for a solid ring $\ur A$ one recovers
the topologically nilpotent and power-bounded elements, and one shows that
Huber's \emph{rings of integral elements} $A^{+}$ correspond exactly to the solid
analytic ring structures on $\ur A$ that are cut out by conditions of the form
``$1-f\sigma$ is invertible on $\iHom(P,-)$''. The match with Huber's classical
theory is, in Scholze's words, striking.

Throughout, ``ring'' means commutative ring, everything happens in light
condensed abelian groups, and we freely use solid abelian groups
$\Solid\subseteq\CondAbl$ (\cref{lec:5,lec:6}) and the relative solid theory over
$\mathbb{Z}[T]$ (\cref{lec:7}).

\subsection{Motivation}

By now we have seen several examples of the same package: a light condensed ring
together with a notion of ``complete'' condensed module over it. In each case the
complete modules form a full subcategory of all condensed modules, cut out by a
completeness condition.

\begin{remark}[The recurring examples]
\label{rmk:8:motivation}
The examples encountered so far are
\begin{align}
\label{eq:8:motivation}
&\bigl(\mathbb{Z},\ \Solid_{\mathbb{Z}}\subseteq\Cond(\mathrm{Ab})\bigr),\notag\\
&\bigl(\mathbb{Z}[T],\ \Solid_{\mathbb{Z}[T]}\subseteq\Mod_{\mathbb{Z}[T]}(\Solid_{\mathbb{Z}})\subseteq\Cond(\mathbb{Z}[T])\bigr),\\
&\bigl(\mathbb{Z}_p,\ \Solid_{\mathbb{Z}_p}\bigr),\qquad
\bigl(\mathbb{Q}_p,\ \Solid_{\mathbb{Q}_p}\bigr),\ \dots\notag
\end{align}
The first is the solid theory (\cref{lec:5}); the second is Clausen's affine line,
where one may either take all $\mathbb{Z}[T]$-modules whose underlying group is
solid, or isolate the stronger $\mathbb{Z}[T]$-solid condition
(\cref{def:7:ZT-solid}), which geometrically distinguishes the affine line from
the closed unit disk. In the last two one simply takes the modules whose underlying
condensed abelian group is solid; over $\mathbb{Q}_p$ (or $\mathbb{C}_p$, or a
$p$-adic Banach algebra) there is no meaningful way to strengthen this. The aim
is to axiomatise this situation.
\end{remark}

\subsection{Huber's definitions}

For comparison we recall Huber's basic objects \cite{Huber_1994}; Huber uses
different names, but the content is the following.

\begin{definition}[Huber]
\label{def:8:huber-ring}
A \emph{Huber ring} is a topological ring $A$ that admits an open subring
$A_{0}\subseteq A$ (a \emph{ring of definition}) and a finitely generated ideal
$I\subseteq A_{0}$ (an \emph{ideal of definition}) such that $A_{0}$ carries the
$I$-adic topology.
\end{definition}

\begin{definition}[Ring of integral elements]
\label{def:8:ring-of-integral-elements}
A \emph{ring of integral elements} in a Huber ring $A$ is an open, integrally
closed subring
\begin{equation*}
A^{+}\ \subseteq\ A^{\circ}\ =\ \{\text{power-bounded elements of }A\}.
\end{equation*}
\end{definition}

\begin{definition}[Huber pair]
\label{def:8:huber-pair}
A \emph{Huber pair} is a pair $(A,A^{+})$ of a Huber ring $A$ and a ring of
integral elements $A^{+}$ of $A$.
\end{definition}

\begin{example}[Huber rings]
\label{ex:8:huber-examples}
\begin{enumerate}
\item Any discrete ring $A$ is Huber: take $A_{0}=A$, $I=(0)$.
\item Any ring $A$ with the $I$-adic topology for a finitely generated ideal $I$
is Huber: take $A_{0}=A$, $I=I$.
\item $\mathbb{Q}_p$ is Huber: take $A_{0}=\mathbb{Z}_p$, $I=(p)$ (which is an
ideal in $\mathbb{Z}_p$, not in $\mathbb{Q}_p$). More generally any
nonarchimedean field $K$ is Huber, with $A_{0}$ the valuation ring and $I$ the
ideal generated by a uniformizer; and any Banach algebra over $\mathbb{Q}_p$ (or
another nonarchimedean local field) is Huber, with $A_{0}$ a unit ball and $I$
generated by a uniformizer.
\end{enumerate}
Thus Huber rings are, roughly, localisations of adic rings of the form
(f.g.-ideal)-adic rings, packaged with some extra technology.
\end{example}

\begin{remark}[Completion; only complete rings]
\label{rmk:8:completion}
The completion $\widehat{A}$ of any Huber ring $A$ is again a Huber ring: if
$A\supseteq A_{0}\supseteq I$ as above, then
$\widehat{A}\supseteq\widehat{A_{0}}=$ the $I$-adic completion of $A_{0}$, an open
subring of $\widehat{A}$. Here ``completion'' is the classical one, allowing all
convergent Cauchy sequences. Since
the adic spaces attached to $A$ and to $\widehat{A}$ are definitionally the same,
non-complete Huber rings play no essential role, and \textbf{henceforth all Huber
rings are assumed complete}: whenever we write a Huber pair, the underlying ring
is complete.
\end{remark}

For a complete Huber ring one has the canonical subsets
\begin{equation}
\label{eq:8:Aoo-Acirc}
A^{\circ\circ}\ =\ \{f\in A\mid f^{n}\to0\ \text{as }n\to\infty\}
\ \subseteq\
A^{\circ}\ =\ \{f\in A\mid (f^{n})_{n\ge0}\ \text{is bounded}\}
\ \subseteq\ A,
\end{equation}
the \emph{topologically nilpotent} and \emph{power-bounded} elements. (Since $A$
is complete, $f^{n}\to0$ is a condition, not extra data.) These are canonical,
whereas $A_{0}$ and $I$ are choices; the two are related by
\begin{equation*}
A^{\circ\circ}\ =\ \varinjlim_{I\subseteq A_{0}\subseteq A}\, I,
\qquad
A^{\circ}\ =\ \varinjlim_{A_{0}\subseteq A}\, A_{0},
\end{equation*}
filtered colimits (unions) over all ideals, resp.\ rings, of definition: the
$A_{0}$'s and $I$'s form filtered systems, and $A^{\circ}$ (resp.\
$A^{\circ\circ}$) is their union. For example $(p)\subseteq\mathbb{Z}_p\subseteq
\mathbb{Q}_p$ realises $A^{\circ\circ}\subseteq A^{\circ}\subseteq A$; the
elements of $I$ are topologically nilpotent, but $I=(p^{2})$ is an equally valid
ideal of definition, so the containment is not an equality for a single choice.

\subsection{Analytic rings}

\begin{notation}[Single-letter convention]
\label{not:8:single-letter}
Huber denotes a Huber pair by the single letter $A$, its two constituents being a
topological ring and a ring of integral elements. We follow this convention for
analytic rings: an analytic ring is denoted by a single letter $A$, and its
underlying light condensed ring is denoted $\ur A$ (Clausen's triangle notation,
\cref{lec:1}).
\end{notation}

\begin{definition}[Analytic ring structure]
\label{def:8:analytic-ring}
Let $\ur A$ be a light condensed ring. An \emph{analytic ring structure} on
$\ur A$ is a full subcategory
\begin{equation*}
\Mod(A)\ \subseteq\ \Cond(\ur A)
\end{equation*}
of the light condensed $\ur A$-modules, that is an abelian subcategory and is
\begin{itemize}
\item stable under all limits, all colimits, and all extensions;
\item stable under all $\Ext^{i}_{\ur A}(N,-)$, i.e.\ $\Ext^{i}_{\ur A}(N,M)\in
\Mod(A)$ for $M\in\Mod(A)$ and $N\in\Cond(\ur A)$ (internal Ext);
\item contains $\ur A$;
\end{itemize}
and such that the inclusion $\Mod(A)\hookrightarrow\Cond(\ur A)$ admits a left
adjoint. We call the objects of $\Mod(A)$ the \emph{complete} modules.
\end{definition}

\begin{remark}[Relation to \cref{lec:1}]
\label{rmk:8:equivalent-L1}
This is equivalent to the definition of analytic ring given by Clausen in
\cref{lec:1} (\cref{def:1:analytic-ring}), but phrased from a slightly more
elementary, abelian-categorical perspective: rather than list the long roster of
good properties of $\Solid_{\mathbb{Z}}$ or $\Mod_{\mathbb{Z}[T]}(\Solid)$ every
time, one takes closure under exactly those operations as the axioms. The clause
that the inclusion admits a left adjoint was, in the lecture, initially omitted
and then restored (see \cref{rmk:8:left-adjoint-in-def}).
\end{remark}

\begin{remark}[Grothendieck category; $D(A)$ need not be $D(\Mod(A))$]
\label{rmk:8:grothendieck-heart}
Gabber asked whether the axioms force $\Mod(A)$ to be a Grothendieck abelian
category, i.e.\ to admit a set of generators; Scholze confirmed this is automatic.
A second question --- whether $\Ext$-groups computed in $\Mod(A)$ agree with those
in the ambient $\Cond(\ur A)$ --- is more delicate. The point is that $D(A)$ will
be defined \emph{not} as the derived category of the abelian category $\Mod(A)$,
but as a full subcategory of $D(\Cond(\ur A))$ (\cref{def:8:derived-cat}), of which
$\Mod(A)$ is the \emph{heart} of a $t$-structure (\cref{prop:8:t-structure}). In
most practical cases the two derived categories coincide, but not in general.
\end{remark}

\subsection{Base change and the symmetric monoidal structure}

\begin{proposition}[Base change]
\label{prop:8:base-change}
The inclusion $\Mod(A)\hookrightarrow\Cond(\ur A)$ has a left adjoint
\begin{equation*}
\Cond(\ur A)\longrightarrow\Mod(A),\qquad M\longmapsto M\otimes_{\ur A}A,
\end{equation*}
the \emph{base change}. Its kernel is a $\otimes$-ideal, and $\Mod(A)$ thereby
acquires a unique symmetric monoidal structure making $-\otimes_{\ur A}A$
symmetric monoidal; explicitly
\begin{equation}
\label{eq:8:tensor-A}
M\otimes_{A}N\ :=\ (M\otimes_{\ur A}N)\otimes_{\ur A}A .
\end{equation}
\end{proposition}

\begin{proof}[Proof sketch]
Existence of the left adjoint is formal (\cref{rmk:8:left-adjoint-in-def}). To see
the kernel is a $\otimes$-ideal, let $M\in\Cond(\ur A)$ with $M\otimes_{\ur A}A=0$
and let $N\in\Cond(\ur A)$ be arbitrary; we must show $(N\otimes_{\ur A}M)\otimes_
{\ur A}A=0$, equivalently that
$\Homop_{\ur A}(N\otimes_{\ur A}M,L)=0$ for all $L\in\Mod(A)$. By tensor--Hom
adjunction
\begin{equation*}
\Homop_{\ur A}(N\otimes_{\ur A}M,L)\ =\
\Homop_{\ur A}\bigl(M,\ \underline{\Homop}_{\ur A}(N,L)\bigr),
\end{equation*}
and $\underline{\Homop}_{\ur A}(N,L)\in\Mod(A)$ by the internal-Ext stability of
\cref{def:8:analytic-ring}; since $M\otimes_{\ur A}A=0$, i.e.\ $M$ has no maps to
any complete module, this vanishes. The tensor product
\eqref{eq:8:tensor-A} is then forced: one checks
\begin{equation*}
(M\otimes_{\ur A}N)\otimes_{\ur A}A\ \xrightarrow{\ \sim\ }\
\bigl((M\otimes_{\ur A}A)\otimes_{\ur A}(N\otimes_{\ur A}A)\bigr)\otimes_{\ur A}A
\end{equation*}
for all $M,N\in\Cond(\ur A)$ by the same argument as in the solid case
(\cref{lec:5}).
\end{proof}

\begin{remark}[Why left-adjointness belongs in the definition]
\label{rmk:8:left-adjoint-in-def}
Scholze first hoped to \emph{prove} the left adjoint exists: there is a general
principle that a full subcategory of a presentable category, closed under all
limits and all sufficiently filtered colimits, is reflective (cf.\
\cref{rmk:6:adjoint-existence}). The non-trivial point is the existence of a generating
set. However, a genuine gap appears one level up, at the derived stage: stability
of $D(A)$ under \emph{arbitrary} (uncountable) products cannot be had for free,
because arbitrary products in $\Cond(\ur A)$ are not exact. The clean fix,
already adopted by Clausen in \cref{lec:1}, is to \emph{require} the left adjoint
to exist as part of \cref{def:8:analytic-ring}; arbitrary products are then stable
because a category admitting a left adjoint to a product-preserving inclusion has
well-behaved products.
\end{remark}

\begin{remark}[The analytic ring structure is an ``induced'' localisation]
\label{rmk:8:localisation}
So an analytic ring structure is a localisation-type situation: $\Mod(A)$ is a
reflective symmetric monoidal subcategory of the condensed modules over the
underlying ring, and the reflector $-\otimes_{\ur A}A$ is a symmetric monoidal
localisation whose kernel is a $\otimes$-ideal. This is the abelian shadow of the
induced analytic ring structures of \cref{rmk:7:analytic-ring}.
\end{remark}

\subsection{Derived categories of analytic rings}

\begin{definition}[Derived category of an analytic ring]
\label{def:8:derived-cat}
Let $A$ be an analytic ring structure on $\ur A$. The \emph{derived category}
$D(A)\subseteq D(\Cond(\ur A))$ is the full subcategory
\begin{equation*}
D(A)\ =\ \bigl\{\, M\in D(\Cond(\ur A))\ \bigm|\ H_{i}(M)\in\Mod(A)\ \text{for all }i\in\mathbb{Z}\,\bigr\}
\end{equation*}
(homological grading). There is a natural comparison functor
$D(\Mod(A))\to D(A)$, but it is \emph{not} always an equivalence (though it is in
essentially all cases arising in practice).
\end{definition}

\begin{proposition}[Structure of $D(A)$]
\label{prop:8:D-triangulated}
$D(A)\subseteq D(\Cond(\ur A))$ is a triangulated subcategory, stable under all
direct sums and all products. The inclusion admits a left adjoint, the
\emph{derived base change}
\begin{equation*}
-\otimes^{\mathbb{L}}_{\ur A}A\colon\ D(\Cond(\ur A))\longrightarrow D(A),
\end{equation*}
whose kernel is a $\otimes$-ideal, and there is a unique symmetric monoidal
tensor product $-\otimes^{\mathbb{L}}_{A}-$ on $D(A)$ making
$-\otimes^{\mathbb{L}}_{\ur A}A$ symmetric monoidal.
\end{proposition}

\begin{proof}[Proof sketch]
\emph{Triangulated.} Given a triangle $M'\to M\to M''\to M'[1]$ in
$D(\Cond(\ur A))$ with $M',M''\in D(A)$, the long exact homology sequence
\begin{equation*}
\cdots\to H_{i+1}(M'')\to H_{i}(M')\to H_{i}(M)\to H_{i}(M'')\to H_{i-1}(M')\to\cdots
\end{equation*}
exhibits each $H_{i}(M)$ as an extension of a subobject $K$ of $H_i(M'')$ by a
quotient $Q$ of $H_i(M')$, both in $\Mod(A)$; stability under kernels, quotients
and extensions gives $H_{i}(M)\in\Mod(A)$, so $M\in D(A)$.

\emph{Direct sums and products.} Stability under $\bigoplus$ reduces to the
$H_{i}$'s. Countable products $\prod_{\mathbb{N}}$ likewise reduce to homology;
arbitrary (uncountable) products are fine only once one admits a left adjoint
(\cref{rmk:8:left-adjoint-in-def}).

\emph{Kernel is a $\otimes$-ideal.} Let $M\in\ker(D(\Cond(\ur A))\to D(A))$ and
$N\in D(\Cond(\ur A))$; we must show $\Homop_{D(\Cond(\ur A))}(M\otimes^{\mathbb{L}}
_{\ur A}N,L)=0$ for all $L\in D(A)$, i.e.\ that $R\iHom(N,L)\in D(A)$. Here light
condensed sets are \emph{replete} (countable limits of surjections are
surjections, \cref{prop:2:countable-limits}), so every $K\in D(\Cond(\ur A))$ is
its own Postnikov limit,
\begin{equation}
\label{eq:8:postnikov}
K\ \xrightarrow{\ \sim\ }\ R\varprojlim_{n}\tau_{\le n}K .
\end{equation}
This lets one reduce (writing $L\in D^{+}$, $N=\operatorname*{colim}_{n}\tau_{\ge-n}N
\in D^{-}$) via a spectral sequence to the abelian statement: for $L\in\Mod(A)$
and $N\in\Cond(\ur A)$, the condition $\Ext^{i}_{\ur A}(N,L)\in\Mod(A)$ --- which
holds by \cref{def:8:analytic-ring}. The boundedness furnished by
\eqref{eq:8:postnikov} is essential: the internal-Ext stability is available only
after such a reduction, since the complexes are a priori unbounded in both
directions.

\emph{Tensor product.} Existence of $-\otimes^{\mathbb{L}}_{A}-$ is then formal,
exactly as in \cref{prop:8:base-change}.
\end{proof}

\begin{proposition}[$t$-structure and heart]
\label{prop:8:t-structure}
$D(A)$ carries a natural $t$-structure, for which the inclusion
$D(A)\hookrightarrow D(\Cond(\ur A))$ is $t$-exact and the derived base change
$-\otimes^{\mathbb{L}}_{\ur A}A$ preserves connective objects $D_{\ge0}$. Its
heart is $\Mod(A)$,
\begin{equation*}
D(A)^{\heartsuit}\ =\ \Mod(A),
\end{equation*}
and taking hearts recovers the abelian-level functors:
$(-\otimes^{\mathbb{L}}_{\ur A}A)^{\heartsuit}=-\otimes_{\ur A}A$ and
$(-\otimes^{\mathbb{L}}_{A}-)^{\heartsuit}=-\otimes_{A}-$.
\end{proposition}

\begin{proof}
$D(A)$ is stable under all $\tau_{\le n}$ and $\tau_{\ge n}$ by
\cref{def:8:derived-cat}, giving the $t$-structure with
$D(A)^{\heartsuit}=\Mod(A)$. A left adjoint to a $t$-exact functor preserves
connective objects, so $-\otimes^{\mathbb{L}}_{\ur A}A$ preserves $D_{\ge0}$. The
heart identities are immediate.
\end{proof}

\begin{remark}[Derived base change is not $t$-exact]
\label{rmk:8:not-t-exact}
The derived base change $-\otimes^{\mathbb{L}}_{\ur A}A$ preserves connective
objects but is not in general $t$-exact --- as already seen for solidification,
which can push a discrete object to the left (\cref{prop:6:singular-homology},
\cref{rmk:7:not-t-exact}). The comparison $D(\Mod(A))\to D(A)$ and the
question of whether animating the abelian tensor product recovers
$-\otimes^{\mathbb{L}}_{A}-$ both fail in general; when the product is recovered,
its underlying abelian functor is automatically the right one, by adjunction.
\end{remark}

\subsection{Maps of analytic rings}

\begin{definition}[Map of analytic rings]
\label{def:8:maps-analytic-rings}
A \emph{map of analytic rings} $A=(\ur A,\Mod(A))\to B=(\ur B,\Mod(B))$ is a map
$\ur A\to\ur B$ of underlying condensed rings such that restriction of scalars
sends complete modules to complete modules, i.e.\ the square
\begin{equation*}
\begin{tikzcd}[column sep=large, row sep=large]
\Mod(B)\arrow[r, hook]\arrow[d, "\mathrm{res}"'] & \Cond(\ur B)\arrow[d, "\mathrm{res}"] \\
\Mod(A)\arrow[r, hook] & \Cond(\ur A)
\end{tikzcd}
\end{equation*}
commutes (the right vertical map being restriction along $\ur A\to\ur B$).
Passing to left adjoints yields the \emph{base change}
\begin{equation*}
-\otimes_{A}B\colon\ \Mod(A)\longrightarrow\Mod(B),
\end{equation*}
a symmetric monoidal ($\otimes$-)functor, together with its derived version
$-\otimes^{\mathbb{L}}_{A}B\colon D(A)\to D(B)$ (which one may compute by base
changing as condensed modules and then completing). The passage to left adjoints
is valid already on hearts, hence on derived categories, since it can be checked
on underlying objects.
\end{definition}

\subsection{Comparison with Huber rings}

Huber rings (assumed complete, \cref{rmk:8:completion}) define condensed rings,
and this functor is fully faithful; in fact it lands in \emph{solid} rings ---
those $\ur A$ whose underlying condensed abelian group is solid. We now recover
Huber's invariants inside the solid formalism.

\begin{definition}[Topologically nilpotent and power-bounded elements]
\label{def:8:tn-pb-solid}
Let $\ur A$ be a solid ring and $f\in\ur A(\ast)$ an element of the underlying
discrete ring, giving a ring map $\mathbb{Z}[T]\to\ur A$, $T\mapsto f$ (the
sequence $1,f,f^{2},\dots$). Following \cref{def:7:tn-pb}, define subsets
$A^{\circ\circ}\subseteq A^{\circ}\subseteq\ur A(\ast)$ by:
\begin{itemize}
\item $f\in A^{\circ\circ}$ (\emph{topologically nilpotent}) if $T\mapsto f$
factors through $\mathbb{Z}[\![T]\!]$,
\begin{equation*}
\begin{tikzcd}[column sep=large, row sep=small]
\mathbb{Z}[T]\arrow[r, "T\mapsto f"]\arrow[d] & \ur A \\
\mathbb{Z}[\![T]\!]\arrow[ur, dashed, "\exists"'] &
\end{tikzcd}
\end{equation*}
i.e.\ the sequence $1,f,f^{2},\dots$ extends to a null sequence (it is enough that
$T\mapsto f$ factors as $\mathbb{Z}[T]$-modules, with the shift relation, cf.\
\cref{rmk:7:factorization-rings});
\item $f\in A^{\circ}$ (\emph{power-bounded}) if $\ur A$ is $\mathbb{Z}[T]$-solid
via $T\mapsto f$, i.e.\ $1-f\sigma$ is an isomorphism on $\iHom(P,\ur A)$, where
$\sigma$ is the shift.
\end{itemize}
When $\ur A$ arises from a Huber ring, these agree with the classical subsets
\eqref{eq:8:Aoo-Acirc}.
\end{definition}

Given a solid ring $\ur A$ equipped with an analytic ring structure $A$, one
extracts a ring of integral elements as follows.

\begin{definition}[Ring of integral elements of an analytic ring]
\label{def:8:Aplus-from-analytic}
Let $A$ be a solid analytic ring structure on $\ur A$. Set
\begin{equation*}
\begin{aligned}
A^{+}=\bigl\{\,f\in\ur A(\ast)\ \bigm|\ & (T\mapsto f)\colon\mathbb{Z}[T]\to\ur A\\
&\text{induces a map of analytic rings }(\mathbb{Z}[T],\Solid_{\mathbb{Z}[T]})\to A\,\bigr\},
\end{aligned}
\end{equation*}
where $(\mathbb{Z}[T],\Solid_{\mathbb{Z}[T]})$ is the closed-unit-disk analytic
ring of \cref{lec:7}. Equivalently, $f\in A^{+}$ iff $1-f\sigma$ is an
automorphism of $P\otimes_{\mathbb{Z}}A$ in $D(A)$ (via the derived base change
$-\otimes^{\mathbb{L}}_{\mathbb{Z}}A$).
\end{definition}

\begin{proposition}[Huber-type structure]
\label{prop:8:Aplus-integrally-closed}
With notation as above,
\begin{equation*}
A^{\circ\circ}\ \subseteq\ A^{+}\ \subseteq\ A^{\circ}\ (\subseteq\ur A(\ast)),
\end{equation*}
$A^{+}$ is an integrally closed subring, and $A^{\circ\circ}$ is a radical ideal.
\end{proposition}

\begin{proof}[Idea]
This is the same argument Clausen gave in \cref{lem:7:Rzero}: everything reduces
to statements of the form ``being solid over one ring implies being solid over
another''. If $f\in A^{\circ\circ}$, then $\ur A$ becomes a
$\mathbb{Z}[\![T]\!]$-module, which is automatically $\mathbb{Z}[T]$-solid, giving
$A^{\circ\circ}\subseteq A^{+}$; the ring/ideal and integral-closedness
statements follow by resolving $\ur A$ by direct sums of the compact projective
generators $\prod_{\mathbb{N}}\mathbb{Z}$ base-changed appropriately, exactly as
in \cref{lec:7}. Indeed the argument there was already phrased for an arbitrary
module, so it applies verbatim.
\end{proof}

Conversely, a ring of integral elements produces an analytic ring structure.

\begin{definition}[Analytic ring of a Huber pair]
\label{def:8:mod-from-huber-pair}
Let $(\ur A,A^{+})$ be a (solid) Huber pair. Define the analytic ring
$(\ur A,A^{+})_{\solid}$ with underlying ring $\ur A$ and complete modules
\begin{equation*}
\begin{aligned}
\Mod\bigl((\ur A,A^{+})_{\solid}\bigr)=\bigl\{\,M\in\Cond(\ur A)\ \bigm|\ &
\forall f\in A^{+},\\
&1-f\sigma\colon\iHom(P,M)\to\iHom(P,M)\ \text{is an isomorphism}\,\bigr\}.
\end{aligned}
\end{equation*}
Declaring $1-f\sigma$ invertible on $\iHom(P,-)$ defines an analytic ring
structure by the general mechanism of \cref{def:8:analytic-ring}; and if
$S\subseteq A^{+}$ generates $A^{+}$ as a ring of integral elements (i.e.\ as an
open integrally closed subring), it suffices to impose the condition for $f\in S$
--- by integral closedness of the resulting power-bounded locus --- so typically
only one or two elements need be checked.
\end{definition}

\begin{example}[{Three structures on $\mathbb{Z}$ and $\mathbb{Z}[T]$}]
\label{ex:8:three-structures}
\begin{align}
\label{eq:8:three-structures}
(\mathbb{Z},\mathbb{Z})_{\solid}&=(\mathbb{Z},\Solid_{\mathbb{Z}})=\Zsolid,\notag\\
(\mathbb{Z}[T],\mathbb{Z})_{\solid}&=\bigl(\mathbb{Z}[T],\Mod_{\mathbb{Z}[T]}(\Solid_{\mathbb{Z}})\bigr),\\
(\mathbb{Z}[T],\mathbb{Z}[T])_{\solid}&=\bigl(\mathbb{Z}[T],\Solid_{\mathbb{Z}[T]}\bigr).\notag
\end{align}
The middle one imposes no condition on $T$ (only solidity over $\mathbb{Z}$),
giving the affine line; the last imposes power-boundedness of $T$, giving the
closed unit disk (\cref{lec:7}).
\end{example}

\begin{proposition}[Huber pairs embed]
\label{prop:8:huber-fully-faithful}
For a Huber pair $(\ur A,A^{+})$, the associated analytic ring
$A=(\ur A,A^{+})_{\solid}$ recovers its ring of integral elements:
\begin{equation*}
\bigl((\ur A,A^{+})_{\solid}\bigr)^{+}\ =\ A^{+}.
\end{equation*}
Consequently the assignment
\begin{equation*}
\{\text{Huber pairs}\}\ \longrightarrow\ \{\text{(solid) analytic rings}\},
\qquad (\ur A,A^{+})\longmapsto(\ur A,A^{+})_{\solid},
\end{equation*}
is fully faithful.
\end{proposition}

\begin{theorem}[Rings of integral elements $\leftrightarrow$ solid analytic structures]
\label{thm:8:correspondence}
Fix a Huber ring $\ur A$. The maps $A\mapsto A^{+}$ and
$A^{+}\mapsto(\ur A,A^{+})_{\solid}$ set up an adjunction
\begin{equation*}
\{\text{rings of integral elements }A^{+}\subseteq\ur A\}\ \rightleftarrows\
\{\text{solid analytic ring structures on }\ur A\},
\end{equation*}
under which $A^{+}\mapsto(\ur A,A^{+})_{\solid}$ is a fully faithful left adjoint.
Thus Huber's rings of integral elements are exactly the analytic ring structures
of the special form ``impose $1-f\sigma$ invertible for $f$ in a subring''.
\end{theorem}

\begin{definition}[Solid analytic ring]
\label{def:8:solid-analytic-ring}
An analytic ring $A=(\ur A,\Mod(A))$ is \emph{solid} if it admits a (necessarily
unique) map $\Zsolid\to A$, where $\Zsolid=(\mathbb{Z},\Solid_{\mathbb{Z}})$.
Equivalently:
\begin{equation*}
A\ \text{solid}\quad\Longleftrightarrow\quad
\text{every }M\in\Mod(A)\ \text{is solid}\quad\Longleftrightarrow\quad
1-\sigma\colon P\otimes_{\mathbb{Z}}A\to P\otimes_{\mathbb{Z}}A\ \text{is an isomorphism}.
\end{equation*}
Uniqueness of the map follows because a map of rings $\mathbb{Z}\to\ur A$ is
unique and solidity is then a condition, not extra structure.
\end{definition}

\begin{remark}[A posteriori motivation]
\label{rmk:8:aposteriori}
Scholze emphasised how closely the solid theory reproduces Huber's classical one:
restricting attention to analytic ring structures cut out by conditions of the
form ``$1-f\sigma$ acts invertibly on $P$'' for $f$ ranging over a subring, one
obtains \emph{precisely} the analytic ring structures induced by Huber's rings of
integral elements. That the ad hoc-looking axiom on $A^{+}$ falls out so exactly
is, he suggested, a strong a posteriori justification for the whole definition of
analytic ring. Feeding a Huber \emph{pair} rather than a bare ring into the
formalism thereby reproduces Huber's theory (cf.\ \cref{rmk:7:analytic-ring}).
\end{remark}
